% \title{Apunte de Matemáticas Discretas CC3101:\\ 
% Soluciones de relaciones de recurrencia  Clase 14}
% % \title{Ejercicios Lógica e Inducción}
% \author{Andrés Abeliuk}
% \date{}
% \maketitle



\section{Recurrencias Lineales}

\subsection{Definición General}

\begin{definitionblock}{Recurrencia lineal homogénea con coeficientes constantes}
Una relación de recurrencia es \textbf{lineal, homogénea y con coeficientes constantes} si tiene la forma:
\[
a_n = c_1 a_{n-1} + c_2 a_{n-2} + \cdots + c_k a_{n-k},
\]
donde $c_1, c_2, \dots, c_k \in \mathbb{R}$ y $c_k \neq 0$. Decimos que la recurrencia es de \textbf{grado $k$}.
\end{definitionblock}

\begin{exampleblock}{Ejemplos}
\begin{itemize}
    \item $a_n = 3a_{n-1} + 2a_{n-3}$ es una recurrencia lineal, homogénea, con coeficientes constantes, de grado $3$.
    \item $a_n = a_{n-1}^2$ \textbf{no} es lineal.
    \item $a_n = a_{n-1} + 1$ \textbf{no} es homogénea.
    \item $a_n = n a_{n-1}$ \textbf{no} tiene coeficientes constantes.
\end{itemize}
\end{exampleblock}

\subsection{Estrategia de Resolución}

La idea central para resolver estas recurrencias es asumir que las soluciones tienen la forma $a_n = r^n$. Sustituyendo esta forma en la ecuación de recurrencia se obtiene la llamada \textbf{ecuación característica}:

\begin{align*}
 &  r^n = c_1 r^{n-1} + c_2 r^{n-2} + \cdots + c_k r^{n-k} \\
& \hspace{-3ex} \Leftrightarrow\\
&  r^k = c_1 r^{k-1} + c_2 r^{k-2} + \cdots + c_k
\end{align*}

Resolver esta ecuación polinomial nos da las raíces $r_i$, y la solución general de la recurrencia se construye a partir de estas raíces, dependiendo de su multiplicidad.

\subsubsection{Caso de Grado 2 con Raíces Distintas}

\begin{thmblock}{Solución con raíces distintas}
Sea la recurrencia:
\[
a_n = c_1 a_{n-1} + c_2 a_{n-2}
\]
y supongamos que la ecuación característica $r^2 = c_1 r + c_2$ tiene dos raíces reales distintas $r_1$ y $r_2$. Entonces, la solución general es:
\[
a_n = \alpha_1 r_1^n + \alpha_2 r_2^n,
\]
donde $\alpha_1$ y $\alpha_2$ se determinan usando las condiciones iniciales.
\end{thmblock}
\begin{proof}
Por hipótesis, sabemos que $r_1^2 = c_1 r_1 + c_2$ y $r_2^2 = c_1 r_2 + c_2$. Entonces:
\begin{align*}
c_1 a_{n-1} + c_2 a_{n-2} &= c_1(\alpha_1 r_1^{n-1} + \alpha_2 r_2^{n-1}) + c_2(\alpha_1 r_1^{n-2} + \alpha_2 r_2^{n-2}) \\
&= \alpha_1 r_1^{n-2}(c_1 r_1 + c_2) + \alpha_2 r_2^{n-2}(c_1 r_2 + c_2) \\
&= \alpha_1 r_1^n + \alpha_2 r_2^n \\
&= a_n
\end{align*}
Concluimos que $a_n = \alpha_1 r_1^n + \alpha_2 r_2^n$ es solución para cualquier par de constantes $\alpha_1, \alpha_2$.

Si además se conocen las condiciones iniciales $a_0$ y $a_1$, entonces las constantes $\alpha_1, \alpha_2$ se determinan unívocamente resolviendo el sistema lineal:
\[
\begin{cases}
a_0 = \alpha_1 r_1^0 + \alpha_2 r_2^0 \\
a_1 = \alpha_1 r_1^1 + \alpha_2 r_2^1
\end{cases}
\]
\end{proof}

\begin{exampleblock}{Ejemplo}
Considere la relación de recurrencia,
\[
\begin{cases}
a_n = 2a_{n-1} + 3a_{n-2}, \quad \text{para } n \geq 2 \\
a_0 = a_1 = 1
\end{cases}
\]
La ecuación característica es $r^2 = 2r + 3 \Rightarrow r^2 - 2r - 3 = 0$, con raíces $r_1 = 3$, $r_2 = -1$.

Solución general:
\[
a_n = \alpha_1 3^n + \alpha_2 (-1)^n
\]
Usamos condiciones iniciales para determinar:
\[
\begin{cases}
1 = \alpha_1 + \alpha_2 \\
1 = 3\alpha_1 - \alpha_2
\end{cases}
\Rightarrow \alpha_1 = \frac{1}{2}, \quad \alpha_2 = \frac{1}{2}
\]
Solución final:
\[
a_n = \tfrac{1}{2} 3^n + \tfrac{1}{2} (-1)^n
\]
\end{exampleblock}

\subsection{Raíces Dobles}

\begin{thmblock}{Solución con raíz doble}
Si la ecuación característica tiene una raíz doble $r_*$, es decir, $r^2 = c_1 r + c_2$ con $r_1 = r_2 = r_*$, entonces la solución general es:
\[
a_n = \alpha_1 r_*^n + \alpha_2 n r_*^n.
\]
\end{thmblock}
\begin{proof}
Por hipótesis, sabemos que $r_*^2 = c_1 r_* + c_2$. Por lo tanto:
\begin{align*}
c_1 a_{n-1} + c_2 a_{n-2} &= c_1(\alpha_1 r_*^{n-1} + \alpha_2 (n-1) r_*^{n-1}) + c_2(\alpha_1 r_*^{n-2} + \alpha_2 (n-2) r_*^{n-2}) \\
&= \alpha_1 r_*^{n-2} (c_1 r_* + c_2) + \alpha_2 r_*^{n-2} ((n-1) c_1 r_* + (n-2) c_2) \\
&= \alpha_1 r_*^n + \alpha_2 r_*^{n-2} (n c_1 r_* - c_1 r_* + n c_2 - 2 c_2) \\
&= \alpha_1 r_*^n + \alpha_2 r_*^{n-2} (n(c_1 r_* + c_2) - c_1 r_* - 2c_2) \\
&= \alpha_1 r_*^n + \alpha_2 n r_*^n - \alpha_2 r_*^{n-2}(c_1 r_* + 2c_2) \\
&= a_n - \alpha_2 r_*^{n-2}(c_1 r_* + 2c_2)
\end{align*}
Para que la expresión sea igual a $a_n$, basta establecer que $c_1 r_* + 2c_2 = 0$.

\noindent \textit{Nota:} Como $r_*$ es raíz doble, esto significa que la ecuación característica
\[
r^2 - c_1 r - c_2 = 0
\]
tiene un discriminante nulo:
\[
\Delta = c_1^2 + 4c_2 = 0.
\]
Aplicando la fórmula general de las raíces cuadráticas:
\[
r = \frac{c_1 \pm \sqrt{c_1^2 + 4c_2}}{2},
\]
obtenemos que:
\[
r_* = \frac{c_1}{2}
\]
cuando el discriminante es cero. Reemplazando en $c_1 r_* + 2c_2$ se obtiene:
\[
c_1 r_* + 2c_2 = \frac{c_1^2}{2} + 2c_2 = 2\Delta =  0,
\]
por lo tanto, se verifica la igualdad y concluimos que
\[
a_n = \alpha_1 r_*^n + \alpha_2 n r_*^n
\]
es solución para cualquier par de constantes $\alpha_1, \alpha_2$.
\end{proof}


\subsection{Caso General: Grado Arbitrario}

\begin{thmblock}{Solución general de grado $k$}
Sea
\[
a_n = c_1 a_{n-1} + \cdots + c_k a_{n-k}
\]
una recurrencia con ecuación característica que tiene $t$ raíces distintas $r_1, \dots, r_t$ con multiplicidades $m_1, \dots, m_t$. Entonces, la solución general es:

\[
a_n = \sum_{i=1}^t \left( \sum_{j=0}^{m_i - 1} \alpha_{i,j} n^j \right) r_i^n
\]
\end{thmblock}

% \subsection*{Ejercicio: Recurrencias simultáneas}
% \noindent Resuelva el siguiente sistema de recurrencias:
% \[
% \begin{cases}
% a_n = 3 a_{n-1} + 2 b_{n-1} \\
% b_n = a_{n-1} + 2 b_{n-1} \\
% a_0 = 1,\ b_0 = 2
% \end{cases}
% \]
% \textbf{Paso 1:} Resolvamos la segunda ecuación. Se puede reescribir como una recurrencia lineal para $b_n$:
% \[
% b_n = a_{n-1} + 2b_{n-1}.
% \]
% Pero aún depende de $a_{n-1}$, así que intentemos eliminar $a_n$ utilizando la diferencia:
% Restamos la primera y la segunda ecuación:
% \[
% a_n - b_n = 2(a_{n-1} - b_{n-1})
% \]
% Definamos $d_n = a_n - b_n$. Entonces:
% \[
% d_n = 2 d_{n-1}, \quad d_0 = a_0 - b_0 = 1 - 2 = -1.
% \]
% Esto nos da:
% \[
% d_n = -1 \cdot 2^n = -2^n.
% \]
% Ahora usamos que $a_n = b_n + d_n = b_n - 2^n$ en la ecuación de $b_n$:
% \[
% b_n = a_{n-1} + 2 b_{n-1} = (b_{n-1} - 2^{n-1}) + 2 b_{n-1} = 3b_{n-1} - 2^{n-1}
% \]
% Entonces $b_n$ cumple:
% \[
% b_n = 3b_{n-1} - 2^{n-1}, \quad b_0 = 2.
% \]
% Esta es una recurrencia no homogénea. Propongamos una solución particular de la forma $A 2^n$. Probamos:
% \[
% A 2^n = 3A 2^{n-1} - 2^{n-1} \Rightarrow A 2^n = (3A - 1) 2^{n-1} \Rightarrow A = 3A - 1 \Rightarrow A = \frac{1}{2}.
% \]
% Solución particular: $b_n^{(p)} = \frac{1}{2} 2^n$.\
% Solución general de la homogénea: $b_n^{(g)} = \alpha 3^n$.\
% Entonces:
% \[
% b_n = \alpha 3^n + \frac{1}{2} 2^n.
% \]
% Con $b_0 = 2$ obtenemos $\alpha = \frac{3}{2}$, así que:
% \[
% b_n = \frac{3}{2} 3^n + \frac{1}{2} 2^n.
% \]
% Finalmente, usamos que $a_n = b_n - 2^n$:
% \[
% a_n = \frac{3}{2} 3^n + \frac{1}{2} 2^n - 2^n = \frac{3}{2} 3^n - \frac{1}{2} 2^n.
% \]
% \textbf{Solución final:}
% \[
% \boxed{a_n = \frac{3}{2} 3^n - \frac{1}{2} 2^n, \quad b_n = \frac{3}{2} 3^n + \frac{1}{2} 2^n.}
% \]



\subsection{Recurrencias No Homogéneas}
\begin{definitionblock}{Recurrencia lineal no homogénea}
Una relación de recurrencia lineal, a coeficientes constantes de grado $k$, se dice \textbf{no homogénea} si es de la forma:
\[
a_n = c_1 a_{n-1} + c_2 a_{n-2} + \cdots + c_k a_{n-k} + f(n),
\]
donde $c_1, c_2, \dots, c_k$ son números reales, $c_k \neq 0$, y $f(n)$ es una función arbitraria. \\
Definimos la \textbf{recurrencia homogénea asociada}:
\[
a_n = c_1 a_{n-1} + c_2 a_{n-2} + \cdots + c_k a_{n-k}
\]
\end{definitionblock}

\begin{exampleblock}{Ejemplos}
\begin{itemize}
  \item $a_n = 2a_{n-1} + 1$ \hfill (no homogénea de grado 1)
  \item $a_n = 2a_{n-1} - 3a_{n-3} + 5n^2 + 2^n$ \hfill (no homogénea de grado 3)
\end{itemize}
\end{exampleblock}

\begin{thmblock}{Teorema de la solución general}
La solución general de una recurrencia no homogénea es:
\[
a_n = a_n^{(g)} + a_n^{(p)},
\]
donde $a_n^{(g)}$ es la solución general de la homogénea asociada, y $a_n^{(p)}$ es una solución particular de la no homogénea.
\end{thmblock}

% \noindent \textit{¿Qué es una solución particular?}\\
\noindent Una \textbf{solución particular} $a_n^{(p)}$ es una función concreta que satisface la relación de recurrencia completa, es decir, incluyendo el término no homogéneo $f(n)$. A diferencia de la solución general, que contiene constantes libres, una solución particular es un caso específico. Se utiliza como complemento a la solución general de la parte homogénea para construir la solución total.\\
Por ejemplo, si tenemos:
\[
a_n = 2a_{n-1} + 1,
\]
la recurrencia homogénea asociada es $a_n^{(g)} = \alpha 2^n$. Una solución particular válida es $a_n^{(p)} = -1$, ya que satisface:
\[
-1 = 2(-1) + 1.
\]
Entonces, la solución general completa es:
\[
a_n = \alpha 2^n - 1.
\]

\begin{proof}
Como $a_n^{(p)}$ es una solución particular de la relación de recurrencia no homogénea, se cumple que:
\[
a^{(p)}_n = c_1a^{(p)}_{n-1} + c_2a^{(p)}_{n-2} + \dots + c_ka^{(p)}_{n-k} + f(n)
\]
Supongamos ahora que $b_n$ es otra solución de la misma relación no homogénea, entonces también se cumple:
\[
b_n = c_1b_{n-1} + c_2b_{n-2} + \dots + c_kb_{n-k} + f(n)
\]
Restando ambas expresiones:
\[
b_n - a^{(p)}_n = c_1(b_{n-1} - a^{(p)}_{n-1}) + c_2(b_{n-2} - a^{(p)}_{n-2}) + \dots + c_k(b_{n-k} - a^{(p)}_{n-k})
\]
Esto implica que la diferencia $b_n - a_n^{(p)}$ satisface la relación de recurrencia homogénea asociada. Es decir, $b_n - a_n^{(p)} = a_n^{(g)}$ para alguna solución homogénea $a_n^{(g)}$.\
Por lo tanto:
\[
b_n = a_n^{(p)} + a_n^{(g)},
\]
como se quería demostrar.
\end{proof}

\begin{exampleblock}{Ejemplo: Torres de Hanoi}
La mínima cantidad de movimientos $h_n$ necesarios para resolver el problema de las Torres de Hanoi con $n$ discos satisface:
\[
\begin{cases}
h_1 = 1 \\
h_n = 2 h_{n-1} + 1 \quad (n \geq 2)
\end{cases}
\]

\textbf{Paso 1:} Se trata de una recurrencia no homogénea de primer orden con $f(n)=1$.

La recurrencia homogénea asociada es:
\[
h_n = 2 h_{n-1},
\]
cuyo comportamiento general es:
\[
h_n^{(g)} = \alpha 2^n.
\]

\textbf{Paso 2:} Buscamos una solución particular $h_n^{(p)}$ tal que:
\[
h_n^{(p)} = 2 h_{n-1}^{(p)} + 1.
\]
Probamos con una constante $h_n^{(p)} = C$:
\[
C = 2C + 1 \Rightarrow C = -1.
\]
Entonces, una solución particular es $h_n^{(p)} = -1$.

\textbf{Paso 3:} La solución general es:
\[
h_n = h_n^{(g)} + h_n^{(p)} = \alpha 2^n - 1.
\]
Usamos la condición inicial $h_1 = 1$:
\[
1 = \alpha 2^1 - 1 \Rightarrow 1 = 2\alpha - 1 \Rightarrow \alpha = 1.
\]

\textbf{Solución cerrada:}
\[
\boxed{h_n = 2^n - 1}
\]
\end{exampleblock}

\subsection{Soluciones Particulares: Método de coeficientes indeterminados}

Obtener la solución general $a_n^{(g)}$ de la recurrencia homogénea asociada es fácil (vimos cómo hacerlo).  Sin embargo no hay ningún mecanismo general para obtener una solución particular $a_n^{(p)}$ de una
recurrencia no homogénea. 
Para una clase particular de $f$'s tenemos el siguiente teorema:


\begin{thmblock}{Teorema}
Sea la recurrencia $a_n = c_1 a_{n-1} + c_2 a_{n-2} + \cdots + c_k a_{n-k} + f(n)$, donde
\[ f(n) = (b_t n^t + b_{t-1} n^{t-1} + \cdots + b_1 n + b_0) s^n, \]
y $b_0, b_1, \dots, b_t$ y $s$ son números reales. Entonces:

\begin{itemize}
  \item Si $s$ \textbf{no es raíz} de la ecuación característica asociada, entonces existe una solución particular de la forma:
  \[ (p_t n^t + p_{t-1} n^{t-1} + \cdots + p_1 n + p_0) s^n. \]

  \item Si $s$ \textbf{es raíz} de la ecuación característica con multiplicidad $m$, entonces existe una solución particular de la forma:
  \[ n^m (p_t n^t + p_{t-1} n^{t-1} + \cdots + p_1 n + p_0) s^n. \]

\end{itemize}
\end{thmblock}

\subsection*{Ejercicio resuelto}

\begin{exampleblock}{Ejercicio: recurrencia con término lineal}
Encuentre la solución general de la recurrencia:
\[
a_n = 3 a_{n-1} - 4n
\]
\end{exampleblock}

\noindent \textbf{Paso 1:} Estudiamos la recurrencia homogénea asociada:
\[
a_n^{(g)} = 3 a_{n-1}^{(g)} \Rightarrow a_n^{(g)} = \alpha 3^n
\]

\noindent \textbf{Paso 2:} Como el término no homogéneo es $f(n) = -4n$, un polinomio de grado 1, proponemos una solución particular de la forma:
\[
a_n^{(p)} = cn + d
\]
Sustituimos en la ecuación original:
\[
cn + d = 3(c(n-1) + d) - 4n
\]
Expandimos el lado derecho:
\[
cn + d = 3cn - 3c + 3d - 4n
\]
Agrupamos términos:
\[
cn + d = (3c - 4)n + (3d - 3c)
\]
Igualando coeficientes de ambos lados:
\[
\begin{cases}
c = 3c - 4 \\
d = 3d - 3c
\end{cases} \Rightarrow 
\begin{cases}
2c = 4 \Rightarrow c = 2 \\
3d - 3c = d \Rightarrow 2d = 3c \Rightarrow d = \frac{3c}{2} = 3
\end{cases}
\]

\noindent \textbf{Paso 3:} Por lo tanto, una solución particular es:
\[
a_n^{(p)} = 2n + 3
\]
\textbf{Solución general:}
\[
a_n = a_n^{(g)} + a_n^{(p)} = \alpha 3^n + 2n + 3
\]
\textbf{Nota:} Para determinar $\alpha$ se requiere una condición inicial adicional, como $a_0$ o $a_1$.


\subsection{Ejercicios Propuestos}

\begin{itemize}
    \item Resolver $a_n = a_{n-1} + 2a_{n-2}$, con $a_0 = 2$, $a_1 = 7$.
    \item Encontrar una fórmula explícita para los números de Fibonacci.
    % \item Hallar la solución de $a_n = 3 a_{n-1} - 4n$.
    \item Encontrar una solución particular para: $a_n = 6 a_{n-1} - 9 a_{n-2} + n 3^n$ y $a_n = 6 a_{n-1} - 9 a_{n-2} + (n^2 + 1) 2^n$.
    \item Defina una recurrencia $s_n$ para la suma de los primeros $n$ naturales, y resuélvala.
    \item Resolver el sistema:
    \[
    \begin{cases}
    a_n = 3a_{n-1} + 2b_{n-1} \\
    b_n = a_{n-1} + 2b_{n-1} \\
    a_0 = 1,\ b_0 = 2
    \end{cases}
    \]
\end{itemize}
